\section*{अध्याय \ref{ch:g2:teeth-and-eating-well} --- दाँत और अच्छा खानपान}

\begin{solution}{exo:g2:teeth-and-eating-well:1}
\omterm{def:g2:teeth-and-eating-well:kinds}{कृंतक} काटते हैं; \omterm{def:g2:teeth-and-eating-well:kinds}{रदनक} फाड़ते हैं; \omterm{def:g2:teeth-and-eating-well:kinds}{चर्वणक} पीसते हैं।
\end{solution}

\begin{solution}{exo:g2:teeth-and-eating-well:2}
\omterm{def:g2:teeth-and-eating-well:kinds}{कृंतक} आगे होते हैं (और उनके बग़ल में कोनों पर \omterm{def:g2:teeth-and-eating-well:kinds}{रदनक}); \omterm{def:g2:teeth-and-eating-well:kinds}{चर्वणक} पीछे होते
हैं।
\end{solution}

\begin{solution}{exo:g2:teeth-and-eating-well:3}
\omterm{def:g2:teeth-and-eating-well:kinds}{कृंतक} कौर काट लेते हैं; \omterm{def:g2:teeth-and-eating-well:kinds}{रदनक} उसे पकड़कर खींचने में मदद करते हैं; जीभ
उसे पीछे ले जाती है; \omterm{def:g2:teeth-and-eating-well:kinds}{चर्वणक} उसे लुगदी बना देते हैं; फिर उसे निगला
जाता है।
\end{solution}

\begin{solution}{exo:g2:teeth-and-eating-well:4}
\omterm{prop:g2:teeth-and-eating-well:sets}{दूध के दाँत} $20$ होते हैं। लगभग छह साल की उम्र से वे ढीले होकर एक-एक
करके गिरते हैं, और नीचे से उगते \omterm{prop:g2:teeth-and-eating-well:sets}{स्थायी दाँत} उन्हें धकेलते हैं।
\end{solution}

\begin{solution}{exo:g2:teeth-and-eating-well:5}
क्योंकि वे आख़िरी सेट हैं: \omterm{def:g2:animals-grow-up:born}{दूध} का गिरा दाँत \omterm{prop:g2:teeth-and-eating-well:sets}{स्थायी दाँत} से बदल जाता
है, पर स्थायी गिरा दाँत कभी नहीं बदलता। स्थायी सेट को पूरी उम्र चलना
है।
\end{solution}

\begin{solution}{exo:g2:teeth-and-eating-well:6}
फल और सब्ज़ियाँ; रोटी और अनाज; \omterm{def:g2:animals-grow-up:born}{दूध} से बनी चीज़ें (\omterm{def:g2:animals-grow-up:born}{दूध}, दही, पनीर);
माँस, मछली या \omterm{def:g2:animals-grow-up:born}{अंडे}; और मीठी चीज़ें सिर्फ़ थोड़ी और कभी-कभार। शरीर को
सचमुच ज़रूरी अकेला पेय पानी है।
\end{solution}

\begin{solution}{exo:g2:teeth-and-eating-well:7}
सेब से विटामिन, रोटी से टिकाऊ ऊर्जा, और पानी --- यानी असली पोषण, और
दाँतों के लिए ईमानदार चबाने का काम। थैला ख ज़्यादातर चीनी देता है, जो
उलटे दाँतों पर हमला करने वाले \omterm{def:g1:healthy-habits:germs}{रोगाणुओं} का पोषण करती है।
\end{solution}

\begin{solution}{exo:g2:teeth-and-eating-well:8}
उत्तर खुला है --- आईने में: आगे चपटे तेज़ धार वाले \omterm{def:g2:teeth-and-eating-well:kinds}{कृंतक}, हर कोने पर
एक नुकीला \omterm{def:g2:teeth-and-eating-well:kinds}{रदनक}, पीछे चौड़े उभरे हुए \omterm{def:g2:teeth-and-eating-well:kinds}{चर्वणक}, ठीक चित्र जैसे।
\end{solution}

\begin{solution}{exo:g2:teeth-and-eating-well:9}
चपटे काटने वाले \omterm{def:g2:teeth-and-eating-well:kinds}{कृंतक} और पीसने वाले \omterm{def:g2:teeth-and-eating-well:kinds}{चर्वणक} \omterm{def:g2:what-animals-eat:kinds}{शाकाहारी} ढंग के औज़ार हैं;
नुकीले \omterm{def:g2:teeth-and-eating-well:kinds}{रदनक} \omterm{def:g2:what-animals-eat:kinds}{मांसाहारी} ढंग के। दोनों होने का मतलब है \omterm{def:g1:plants-around-us:plant}{पौधों} और माँस,
दोनों का मेनू: यानी \omterm{def:g2:what-animals-eat:kinds}{सर्वाहारी} का मुँह, उस अध्याय की कलचिड़ी और जंगली
सूअर के मेनू जैसा।
\end{solution}

\begin{solution}{exo:g2:teeth-and-eating-well:10}
छूट रहा है: फल और सब्ज़ियाँ (पूरे दिन एक भी नहीं!) और रोज़ के पेय के
रूप में पानी; माँस, मछली या \omterm{def:g2:animals-grow-up:born}{अंडे} भी ग़ायब हैं। बहुत ज़्यादा बार: मीठी
चीज़ें --- मिठाई, ठंडा पेय, बिस्कुट। एक सुधार: नाश्ते और शाम के नाश्ते
में फल, दोपहर या रात के खाने में सब्ज़ियाँ और थोड़ी मछली या \omterm{def:g2:animals-grow-up:born}{अंडा}, ठंडे
पेय की जगह पानी, और मीठी चीज़ें दिन में एक ही बार।
\end{solution}
